\subsection{Differentials and derivations}

Let \(\mathbf{B} = \mathbf{A}[(\mathbf{X}_i)_{i\in I}]\) , then by III, p.~569 there exists for each \(i\in I\) one and only one A-derivation \(\mathbf{D}_i\) of \(\mathbf{B}\) such that

\[
\mathrm{D} _ {i} \mathrm{X} _ {i} = 1, \quad \mathrm{D} _ {i} \mathrm{X} _ {j} = 0 \quad \text { for } \quad j \neq i.\tag{5}
\]

The polynomial \(D_{i}P\) is called the partial derivative of P with respect to \(X_{i}\) ; we shall also denote it by \(D_{X_{i}}P\) or \(\frac{\partial P}{\partial X_{i}}\) or \(P'_{X_{i}}\) . By III, p.~558, formula (21), we have, for \(\nu = (\nu_{j}) \in \mathbf{N}^{(I)}\) ,

\[
\mathrm{D} _ {i} \left(\mathbf {X} ^ {\nu}\right) = \left\{ \begin{array}{c c c} \nu_ {i} \mathbf {X} _ {i} ^ {\nu_ {i} - 1} \prod_ {j \in I - \{i \}} \mathbf {X} _ {j} ^ {\nu_ {j}} & \text {if} & \nu_ {i} > 0 \\ 0 & \text {if} & \nu_ {i} = 0. \end{array} \right.\tag{6}
\]

It follows from (6) that \(\mathbf{D}_i\mathbf{D}_j = \mathbf{D}_j\mathbf{D}_i\) for any \(i, j \in I\) . For \(\nu = (\nu_i)_{i \in I} \in \mathbf{N}^{(I)}\) we shall put \(\mathbf{D}^{\nu} = \prod_{i \in I} \mathbf{D}_i^{\nu_i}\) and \(\nu! = \prod_{i \in I} (\nu_i!)\) . With the product ordering on \(\mathbf{N}^{(I)}\) we have

\[
\mathrm{D} ^ {\nu} \left(\mathrm{X} ^ {\mu}\right) = \left\{ \begin{array}{c l} \frac {\mu !}{(\mu - \nu) !} \mathrm{X} ^ {\mu - \nu} & \text { if } \quad \nu \leqslant \mu , \\ 0 & \text { if   not }. \end{array} \right.
\]

When P is a polynomial in a single indeterminate X, the unique partial derivative of P is written DP or \(\frac{dP}{dX}\) or \(P'\) and is called simply the derivative of P.

Again let \(\mathbf{B} = \mathbf{A}\left[(\mathbf{X}_i)_{i\in I}\right]\) ; by III, p.~569 the B-module of A-differentials of B, \(\Omega_{\mathbf{A}}(\mathbf{B})\) , has the family \((d\mathbf{X}_i)_{i\in I}\) of differentials of the \(\mathbf{X}_i\) as basis. Let \(\partial_i\) be the coordinate form of index \(i\) relative to this basis over \(\Omega_{\mathbf{A}}(\mathbf{B})\) . Then the mapping \(u\mapsto \langle \partial_i,du\rangle\) of B into itself is a derivation of B which maps \(\mathbf{X}_i\) to 1 and \(\mathbf{X}_j\) to 0 for \(j\neq i\) , and hence is \(\mathbf{D}_i\) ; in other words, we have

\[
d u = \sum_ {i \in I} \left(\mathrm{D} _ {i} u\right) d \mathrm{X} _ {i}\tag{7}
\]

for each \(u \in \mathbf{B}\) . If \(\mathbf{I}\) is finite, \((\mathbf{D}_i)_{i \in \mathbf{I}}\) is a basis of the B-module of derivations of \(\mathbf{B}\) .

PROPOSITION 4. --- Let E be an associative, commutative and unital A-algebra, \(\mathbf{x} = (x_i)_{i\in I}\) a family of elements of E, \(u\) an element of \(\mathbf{A}[(\mathbf{X}_i)_{i\in I}]\) and \(y = u(\mathbf{x})\) . Then for every derivation D of E into an E-module we have

\[
\mathrm{D} y = \sum_ {i \in I} \left(\mathrm{D} _ {i} u\right) (\mathbf {x}) \cdot \mathrm{D} x _ {i}.
\]

It suffices to prove the proposition when \(u\) is a monomial, and in that case it follows from III, p.~558, Prop. 6.

COROLLARY. --- Let \(f \in \mathbf{A}[\mathbf{X}_1, \ldots, \mathbf{X}_p]\) and \(g_i \in \mathbf{A}[\mathbf{Y}_1, \ldots, \mathbf{Y}_q]\) for \(1 \leqslant i \leqslant p\) , and write \(h = f(g_1, \ldots, g_p)\) , then for \(1 \leqslant j \leqslant q\) we have

\[
\frac {\partial h}{\partial \mathbf {Y} _ {j}} = \sum_ {i = 1} ^ {p} \mathrm{D} _ {i} f (g _ {1}, \dots , g _ {p}) \cdot \frac {\partial g _ {i}}{\partial \mathbf {Y} _ {j}}.\tag{8}
\]

This is the special case \(\mathbf{E} = \mathbf{A}[\mathbf{Y}_1,\dots ,\mathbf{Y}_q]\) , \(x_{i} = g_{i}\) and \(\mathbf{D} = \partial /\partial \mathbf{Y}_j\) of Prop. 4.

Let \(\mathbf{X} = (\mathbf{X}_i)_{i\in I}\) , \(\mathbf{Y} = (\mathbf{Y}_i)_{i\in I}\) be two disjoint families of indeterminates, and write \(\mathbf{X} + \mathbf{Y}\) for the family \((\mathbf{X}_i + \mathbf{Y}_i)_{i\in I}\) . Given \(u\in \mathbf{A}[\mathbf{X}]\) , consider the element \(u(\mathbf{X} + \mathbf{Y})\) of \(\mathbf{A}[\mathbf{X},\mathbf{Y}]\) . For \(\nu \in \mathbf{N}^{(I)}\) we denote by \(\Delta^{\nu}u\) the coefficient of \(\mathbf{Y}^{\nu}\) in \(u(\mathbf{X} + \mathbf{Y})\) , considered as polynomial in the \(\mathbf{Y}_i\) with coefficients in \(\mathbf{A}[\mathbf{X}]\) . By definition we have \(\Delta^{\nu}u\in \mathbf{A}[\mathbf{X}]\) and

\[
u (\mathbf {X} + \mathbf {Y}) = \sum_ {\nu} \left(\Delta^ {\nu} u\right) (\mathbf {X}) \mathbf {Y} ^ {\nu}.\tag{9}
\]

(Here and in the rest of this No., the summations are extended over the index set \(\mathbf{N}^{(I)}\) unless the contrary is said.)

Let \(\mathbf{a} \in \mathbf{A}^{I}\) , then on substituting \(\mathbf{a}\) for \(\mathbf{X}\) and \(\mathbf{X} - \mathbf{a}\) for \(\mathbf{Y}\) in (9) we obtain

\[
u (\mathbf {X}) = \sum_ {\nu} \left(\Delta^ {\nu} u\right) (\mathbf {a}) (\mathbf {X} - \mathbf {a}) ^ {\nu}.\tag{10}
\]

In particular we have

\[
u (\mathbf {X}) = \sum_ {\nu} \left(\Delta^ {\nu} u\right) (0) \mathbf {X} ^ {\nu}.\tag{11}
\]

If \(u, v \in \mathbf{A}[\mathbf{X}]\) , we have

\[
\begin{array}{r l} (u v) (\mathbf {X} + \mathbf {Y}) & = \left(\sum_ {\nu} (\Delta^ {\nu} u) (\mathbf {X}) \mathbf {Y} ^ {\nu}\right) \left(\sum_ {\rho} (\Delta^ {\rho} v) (\mathbf {X}) \mathbf {Y} ^ {\rho}\right) \\ & = \sum_ {\sigma} \left[ \sum_ {\nu + \rho = \sigma} (\Delta^ {\nu} u) (\mathbf {X}) (\Delta^ {\rho} v) (\mathbf {X}) \right] \mathbf {Y} ^ {\sigma} \end{array}
\]

hence

\begin{enumerate}
\def\labelenumi{(\arabic{enumi})}
\setcounter{enumi}{11}
\tightlist
\item
\end{enumerate}

\[
\Delta^ {\sigma} (u v) = \sum_ {\nu + \rho = \sigma} (\Delta^ {\nu} u) (\Delta^ {\rho} v).
\]

Let \(\mathbf{Z} = (Z_i)_{i\in I}\) be another family of indeterminates.

We have :

\[
\begin{array}{r l} \sum_ {\nu} (\Delta^ {\nu} u) (\mathbf {X}) (\mathbf {Y} + \mathbf {Z}) ^ {\nu} & = u (\mathbf {X} + \mathbf {Y} + \mathbf {Z}) \\ & = \sum_ {\sigma} (\Delta^ {\sigma} u) (\mathbf {X} + \mathbf {Y}) \mathbf {Z} ^ {\sigma} \\ & = \sum_ {\rho , \sigma} (\Delta^ {\rho} \Delta^ {\sigma} u) (\mathbf {X}) \mathbf {Y} ^ {\rho} \mathbf {Z} ^ {\sigma}, \end{array}
\]

hence by I, p.~99, Cor. 2 :

\[
\Delta^ {\rho} \Delta^ {\sigma} u = \frac {(\rho + \sigma) !}{\rho ! \sigma !} \Delta^ {\rho + \sigma} u.\tag{13}
\]

PROPOSITION 5. --- For any \(u \in \mathbf{A}[\mathbf{X}]\) and \(\nu \in \mathbf{N}^{(I)}\) we have

\[
\mathbf {D} ^ {\nu} \boldsymbol {u} = \nu ! \Delta^ {\nu} \boldsymbol {u}.
\]

Suppose first that \(\nu\) has length 1; then there exists an element \(i\) of I such that \(\nu = \varepsilon_i\) , that is, \(\nu_i = 1\) and \(\nu_j = 0\) for all \(j \neq i\) in I. Formula (12) shows that \(\Delta^{\varepsilon_i}\) is a derivation of the A-algebra A{[}X{]} which clearly has the value zero on \(\mathbf{X}_j\) for \(j \neq i\) and the value 1 on \(\mathbf{X}_i\) . We thus have \(\Delta^{\varepsilon_i} = \mathbf{D}_i\) for each \(i \in I\) .

By (13) we have

\[
(\rho ! \Delta^ {\rho}) \cdot (\sigma ! \Delta^ {\sigma}) = (\rho + \sigma)! \Delta^ {\rho + \sigma}\tag{14}
\]

in the endomorphism algebra of the A-module \(\mathbf{A}[\mathbf{X}]\) . It now follows by induction on the length of \(\nu\) that \(\nu! \Delta^{\nu} = \mathbf{D}^{\nu}\) .

If \(\mathbf{A}\) is a \(\mathbf{Q}\) -algebra, the formulae (9), (10), (11) may thus be written

\begin{enumerate}
\def\labelenumi{(\arabic{enumi})}
\setcounter{enumi}{14}
\tightlist
\item
\end{enumerate}

\[
u (\mathbf {X} + \mathbf {Y}) = \sum_ {\nu} \frac {1}{\nu !} (\mathrm{D} ^ {\nu} u) (\mathbf {X}) \mathbf {Y} ^ {\nu}
\]

\[
u (\mathbf {X}) = \sum_ {\nu} \frac {1}{\nu !} (D ^ {\nu} u) (\mathbf {a}) (\mathbf {X} - \mathbf {a}) ^ {\nu}\tag{16}
\]

\[
u (\mathbf {X}) = \sum_ {\nu} \frac {1}{\nu !} (D ^ {\nu} u) (0) \mathbf {X} ^ {\nu}.\tag{17}
\]

The formulae (15), (16), (17) are all three called « Taylor's formula ».

PROPOSITION 6 (« Euler's identity »). --- Let \(u \in \mathbf{A}[\mathbf{X}]\) be a homogeneous polynomial of degree \(r\) ; we have

\[
\sum_ {i \in I} X _ {i} \cdot D _ {i} u = r u.
\]

Let D be the A-linear mapping of A{[}X{]} into itself such that \(D(v) = sv\) when v is homogeneous of degree s. We know (III, p.~554, Example 6) that D is a derivation of A{[}X{]}. Thus Prop. 6 is a Corollary of Prop. 4 (IV, p.~6).

